
\begin{lemma}{7.4.90}\label{lem:7-4-90}
若 $L_\infty$ 同态
\[
\widehat{\overline{\mathfrak g}}^{\circ}:
\operatorname{Der}B(\overline C[1])
\longrightarrow
\operatorname{Der}B(\overline C'[1])
\]
诱导同构
\begin{equation}
H\bigl(
\operatorname{gr}(\operatorname{Der}B(\overline C[1])),
\overline\delta
\bigr)
\longrightarrow
H\bigl(
\operatorname{gr}(\operatorname{Der}B(\overline C'[1])),
\overline\delta
\bigr),
\tag{7.4.91}\label{eq:7-4-91}
\end{equation}
则
\begin{equation}
\widehat{\overline{\mathfrak g}}^{\circ}:
H\bigl(
\operatorname{Der}B(\overline C[1]),
\overline\delta
\bigr)
\longrightarrow
H\bigl(
\operatorname{Der}B(\overline C'[1]),
\overline\delta
\bigr)
\tag{7.4.92}\label{eq:7-4-92}
\end{equation}
是同构。
\end{lemma}

注意，由定理 A3.16，
引理~\ref{lem:7-4-90} 的结论蕴含
$\widehat{\overline{\mathfrak g}}^{\circ}$
是同伦等价。

\begin{proof}
在本引理的证明中\translatornote{源书写“在引理 5.1.90 的证明中”；全书并无该编号，而本段正是引理 7.4.90 的证明，故据上下文改写。见疑误 A-0577。}，
用 $\overline Y_k$ 等表示
\[
\operatorname{Hom}\bigl(
B_k(\overline C[1]),
\overline C[1]
\bigr)
\]
中的元素，
用 $\overline y_k$ 等表示
\[
\operatorname{Hom}\bigl(
B_k(\overline C'[1]),
\overline C'[1]
\bigr)
\]
中的元素。
为证明式~\eqref{eq:7-4-92} 的满射性，
使用下面的子引理。

\begin{statement}{子引理}{7.4.93}\label{sublem:7-4-93}\itshape
设
\[
\overline y_i\in
\operatorname{Hom}\bigl(
B_i(\overline C'[1]),
\overline C'[1]
\bigr),
\qquad
\delta(\overline y_1,\ldots,\overline y_i,\ldots)=0.
\]
则对 $i=1,\ldots,k$，
存在
\[
\begin{aligned}
\overline Y_i&\in
\operatorname{Hom}\bigl(
B_i(\overline C[1]),
\overline C[1]
\bigr),\\
\overline z_i&\in
\operatorname{Hom}\bigl(
B_i(\overline C'[1]),
\overline C'[1]
\bigr),
\end{aligned}
\]
使得
\begin{equation}
\delta(\overline Y_1,\ldots,\overline Y_k)
\equiv0
\pmod{F^{k+1}\operatorname{Der}B(\overline C[1])}
\tag{7.4.94}\label{eq:7-4-94}
\end{equation}
以及
\begin{equation}
\begin{aligned}
&\widehat{\overline{\mathfrak g}}^{\circ}
(\overline Y_1,\ldots,\overline Y_k)
-(\overline y_1,\ldots,\overline y_k)\\
&\qquad
-\delta(\overline z_1,\ldots,\overline z_k)
\in F^{k+1}\operatorname{Der}B(\overline C'[1]).
\end{aligned}
\tag{7.4.95}\label{eq:7-4-95}
\end{equation}
此外，若已经给定
$i=1,\ldots,k$ 时的
$\overline Y_i,\overline z_i$，
使式~\eqref{eq:7-4-94}、\eqref{eq:7-4-95}
对 $k$ 成立，
则可选择
$\overline Y_{k+1},\overline z_{k+1}$，
使这两式对 $k+1$ 成立。
\end{statement}

\par\noindent\textit{子引理~\ref{sublem:7-4-93} 的证明.}\quad
对 $k$ 作归纳。
$k=1$ 的情形容易。
假设已有 $i=1,\ldots,k$ 时的
$\overline Y_i,\overline z_i$，
且式~\eqref{eq:7-4-94}、\eqref{eq:7-4-95}
对 $k$ 成立。
置
\begin{equation}
\overline C_{k+1}
\equiv
\delta(\overline Y_1,\ldots,\overline Y_k)
\pmod{F^{k+2}\operatorname{Der}B(\overline C[1])}.
\tag{7.4.96}\label{eq:7-4-96}
\end{equation}
式~\eqref{eq:7-4-94} 对 $k$ 成立蕴含
\[
\overline C_{k+1}\in
\operatorname{Hom}\bigl(
B_{k+1}(\overline C[1]),
\overline C[1]
\bigr).
\]
令
\[
\delta_1:
\operatorname{Hom}\bigl(
B_{k+1}(\overline C[1]),
\overline C[1]
\bigr)
\longrightarrow
\operatorname{Hom}\bigl(
B_{k+1}(\overline C[1]),
\overline C[1]
\bigr)
\]
为由 $\overline{\mathfrak m}_1$ 诱导的余导子。
于是式~\eqref{eq:7-4-96} 蕴含
$\delta_1(\overline C_{k+1})=0$。
由式~\eqref{eq:7-4-95}，
存在
\[
\overline w_{k+1}\in
\operatorname{Hom}\bigl(
B_{k+1}(\overline C'[1]),
\overline C'[1]
\bigr)
\]
使得
\[
\begin{aligned}
&\widehat{\overline{\mathfrak g}}^{\circ}
(\overline Y_1,\ldots,\overline Y_k,0)
-(\overline y_1,\ldots,\overline y_k,\overline y_{k+1})\\
&\quad
-\delta(\overline z_1,\ldots,\overline z_k,0)
\equiv\overline w_{k+1}
\pmod{F^{k+2}\operatorname{Der}B(\overline C'[1])}.
\end{aligned}
\]
这里把模数写成 $F^{k+2}$\translatornote{源书写模 $F^{k+1}$；但等式两端依 (7.4.95) 本就属于 $F^{k+1}$，该读法不能确定第 $k+1$ 元分量，也不能推出紧接着模 $F^{k+2}$ 的计算。故按归纳所需改为 $F^{k+2}$。见疑误 A-0578。}。
计算得
\[
\begin{aligned}
\overline{\mathfrak g}^{\circ}_1(\overline C_{k+1})
&=
\delta\Bigl(
\widehat{\overline{\mathfrak g}}^{\circ}
(\overline Y_1,\ldots,\overline Y_k,0)
\Bigr)\\
&\equiv
\delta(\overline y_1,\ldots,\overline y_k,\overline y_{k+1})
+\delta_1(\overline w_{k+1})\\
&\equiv
\delta_1(\overline w_{k+1})
\pmod{F^{k+2}\operatorname{Der}B(\overline C'[1])}.
\end{aligned}
\]
最后一个同余来自
$\delta(\overline y_1,\ldots,\overline y_i,\ldots)=0$。

因此，由式~\eqref{eq:7-4-91} 是同构，
可找到
\[
\overline Y_{k+1}\in
\operatorname{Hom}\bigl(
B_{k+1}(\overline C[1]),
\overline C[1]
\bigr)
\]
使
\[
\delta_1(\overline Y_{k+1})+\overline C_{k+1}=0.
\]
此外，再次利用式~\eqref{eq:7-4-91} 是同构，
还可选择 $\overline Y_{k+1}$，
使
\[
\overline{\mathfrak g}^{\circ}_1(\overline Y_{k+1})
+\overline w_{k+1}
\]
是一个 $\delta_1$-上边界\translatornote{源书在 $\overline{\mathfrak g}^{\circ}_1(\overline Y_{k+1})$ 后多一枚右括号，并在 $\overline w_{k+1}$ 前写减号。附录定义 A3.12 表明一元分量与微分交换；紧前两式及 (7.4.95) 的下一阶分量都要求这里取加号。译稿据此改正。见疑误 A-0579。}。
故存在 $\overline z_{k+1}$，
其边界等于这个元素。
子引理~\ref{sublem:7-4-93} 得证。\hfill$\square$

子引理~\ref{sublem:7-4-93}
立即蕴含式~\eqref{eq:7-4-92} 的满射性。
为证明其单射性，使用下面的子引理。

\begin{statement}{子引理}{7.4.97}\label{sublem:7-4-97}\itshape
设
\[
\overline Y_i\in
\operatorname{Hom}\bigl(
B_i(\overline C[1]),
\overline C[1]
\bigr)
\]
且
\begin{equation}
\delta(\overline Y_1,\ldots,\overline Y_i,\ldots)=0.
\tag{7.4.98}\label{eq:7-4-98}
\end{equation}
又设
\[
\overline z_i\in
\operatorname{Hom}\bigl(
B_i(\overline C'[1]),
\overline C'[1]
\bigr)
\]
且
\begin{equation}
\widehat{\overline{\mathfrak g}}^{\circ}
(\overline Y_1,\ldots,\overline Y_i,\ldots)
=
\delta(\overline z_1,\ldots,\overline z_i,\ldots).
\tag{7.4.99}\label{eq:7-4-99}
\end{equation}
则对 $i=1,\ldots,k$，存在
\[
\begin{aligned}
\overline Z_i&\in
\operatorname{Hom}\bigl(
B_i(\overline C[1]),
\overline C[1]
\bigr),\\
\overline w_i&\in
\operatorname{Hom}\bigl(
B_i(\overline C'[1]),
\overline C'[1]
\bigr),
\end{aligned}
\]
使得
\begin{equation}
(\overline Y_1,\ldots,\overline Y_k)
\equiv
\delta(\overline Z_1,\ldots,\overline Z_k)
\pmod{F^{k+1}\operatorname{Der}B(\overline C[1])}
\tag{7.4.100}\label{eq:7-4-100}
\end{equation}
以及
\begin{equation}
\begin{aligned}
&\widehat{\overline{\mathfrak g}}^{\circ}
(\overline Z_1,\ldots,\overline Z_k)\\
&\qquad\equiv
(\overline z_1,\ldots,\overline z_k)
+\delta(\overline w_1,\ldots,\overline w_k)
\pmod{F^{k+1}\operatorname{Der}B(\overline C'[1])}.
\end{aligned}
\tag{7.4.101}\label{eq:7-4-101}
\end{equation}
此外，若已经给定
$i=1,\ldots,k$ 时的
$\overline Z_i,\overline w_i$，
使式~\eqref{eq:7-4-100}、\eqref{eq:7-4-101}
对 $k$ 成立，
则可选择
$\overline Z_{k+1},\overline w_{k+1}$，
使这两式对 $k+1$ 成立。
\end{statement}

\par\noindent\textit{子引理~\ref{sublem:7-4-97} 的证明.}\quad
对 $k$ 作归纳。
$k=1$ 的情形容易。
假设已有 $i=1,\ldots,k$ 时的
$\overline Z_i,\overline w_i$，
且式~\eqref{eq:7-4-100}、\eqref{eq:7-4-101}
对 $k$ 成立。
由式~\eqref{eq:7-4-100}，
存在
\[
\overline C_{k+1}\in
\operatorname{Hom}\bigl(
B_{k+1}(\overline C[1]),
\overline C[1]
\bigr)
\]
使得
\begin{equation}
\begin{aligned}
&(\overline Y_1,\ldots,\overline Y_k,\overline Y_{k+1})
-\delta(\overline Z_1,\ldots,\overline Z_k,0)\\
&\qquad\equiv
\overline C_{k+1}
\pmod{F^{k+2}\operatorname{Der}B(\overline C[1])}.
\end{aligned}
\tag{7.4.102}\label{eq:7-4-102}
\end{equation}
于是式~\eqref{eq:7-4-98} 蕴含
\begin{equation}
\delta_1(\overline C_{k+1})=0,
\tag{7.4.103}\label{eq:7-4-103}
\end{equation}
而式~\eqref{eq:7-4-99} 蕴含
\begin{equation}
\begin{aligned}
\overline{\mathfrak g}^{\circ}_1(\overline C_{k+1})
&=
\widehat{\overline{\mathfrak g}}^{\circ}
(\overline Y_1,\ldots,\overline Y_{k+1})\\
&\quad-
\delta\Bigl(
\widehat{\overline{\mathfrak g}}^{\circ}
(\overline Z_1,\ldots,\overline Z_k,0)
\Bigr)\\
&\equiv
\delta(\overline z_1,\ldots,\overline z_{k+1})\\
&\quad-
\delta\Bigl(
\widehat{\overline{\mathfrak g}}^{\circ}
(\overline Z_1,\ldots,\overline Z_k,0)
\Bigr)\\
&\hspace{15mm}
\pmod{F^{k+2}\operatorname{Der}B(\overline C'[1])}.
\end{aligned}
\tag{7.4.104}\label{eq:7-4-104}
\end{equation}
由式~\eqref{eq:7-4-101}，
存在
\[
\overline x_{k+1}\in
\operatorname{Hom}\bigl(
B_{k+1}(\overline C'[1]),
\overline C'[1]
\bigr)
\]
使得
\[
\begin{aligned}
\overline x_{k+1}
&\equiv
\widehat{\overline{\mathfrak g}}^{\circ}
(\overline Z_1,\ldots,\overline Z_k,0)\\
&\quad-
(\overline z_1,\ldots,\overline z_k,\overline z_{k+1})
-\delta(\overline w_1,\ldots,\overline w_k,0)\\
&\hspace{15mm}
\pmod{F^{k+2}\operatorname{Der}B(\overline C'[1])}.
\end{aligned}
\]
式~\eqref{eq:7-4-104} 因而蕴含
\begin{equation}
\overline{\mathfrak g}^{\circ}_1(\overline C_{k+1})
=-\delta_1(\overline x_{k+1}).
\tag{7.4.105}\label{eq:7-4-105}
\end{equation}
所以，由式~\eqref{eq:7-4-91} 是同构，
存在
\[
\begin{aligned}
\overline Z_{k+1}&\in
\operatorname{Hom}\bigl(
B_{k+1}(\overline C[1]),
\overline C[1]
\bigr),\\
\overline w_{k+1}&\in
\operatorname{Hom}\bigl(
B_{k+1}(\overline C'[1]),
\overline C'[1]
\bigr)
\end{aligned}
\]
使得
\[
\overline C_{k+1}=\delta_1(\overline Z_{k+1})
\]
以及
\[
\overline{\mathfrak g}^{\circ}_1(\overline Z_{k+1})
+\overline x_{k+1}
=\delta_1(\overline w_{k+1}).
\]
容易看出，
$\overline Z_{k+1}$ 与 $\overline w_{k+1}$
具有所需性质。
子引理~\ref{sublem:7-4-97} 得证。\hfill$\square$

式~\eqref{eq:7-4-92} 的单射性
立即由子引理~\ref{sublem:7-4-97} 得到。
引理~\ref{lem:7-4-90} 得证。
\end{proof}
